% Round 17. Foundation branch; candidate pending independent review.
\section{Parent retractions on a forest remove all boundary comparisons}
\label{r17fr:section}

\paragraph{Status.}
\textbf{Internally reviewed.} The final claims register identifies the independent audits and frozen dependencies.
This is a Chow-level strengthening of the rooted numerical-surjectivity
argument in FT17.  It does not assume numerical-to-homological
comparison on the double strata.  The retractions correct the
representatives themselves.

Let $\mathcal X/R$ be regular, integral, flat and projective,
with smooth geometrically integral projective generic fiber $X$
of dimension $n\ge3$ over a finite extension of $\mathbf Q_p$.
Suppose the special fiber is reduced and strictly semistable.
Its native components and native connected double strata are
denoted $D_v$ and $E_e$.  They are smooth projective over the
finite residue field $k$, without a geometric-integrality
requirement.  Suppose the resulting finite multigraph is a
forest and orient each tree away from a chosen root.  A triple
intersection would give a triangle, so a forest has no triple
strata.  All Chow and numerical groups have rational coefficients.

\begin{definition}[Parent retraction data]
\label{r17fr:definition}
For each nonroot vertex $v$, let $e$ be its parent edge and
let $j_{e,v}:E_e\hookrightarrow D_v$ be the inclusion.
Require a $k$-morphism
\[
               \pi_v:D_v\longrightarrow E_e,
                 \qquad\pi_vj_{e,v}=\operatorname{id}_{E_e}.
 \tag{FR17.1}\label{r17fr:retraction}
\]
There is no required retraction from the parent component,
no fixed variety common to the different edges, and no
condition on restrictions to the other incident edges.
The root component is unrestricted.
\end{definition}

Put
\[
 C^a=\bigoplus_vN^a(D_v)_{\mathbf Q},\qquad
 W^a=\bigoplus_eN^a(E_e)_{\mathbf Q},\qquad
 (\rho_au)_e=j_{e,v}^*u_v-j_{e,p}^*u_p,
\]
where $p$ and $v$ are the parent and child of $e$.  As before,
$A_a u=(i_v^!\sum_wi_{w*}u_w)_v$.  All degrees use $\deg_k$.

\begin{lemma}[Simultaneous numerical surjectivity]
\label{r17fr:rho-surjective}
Under \eqref{r17fr:retraction}, the map
$\rho_a:C^a\to W^a$ is surjective for every $0\le a\le n-1$.
Moreover
\[
                          \ker A_a=\ker\rho_a.
 \tag{FR17.2}\label{r17fr:kernel}
\]
\end{lemma}

\begin{proof}
Every map is a morphism of smooth projective varieties, so its
Chow pullback preserves numerical equivalence by projection
formula.  In particular $j_{e,v}^*\pi_v^*=\operatorname{id}$
on numerical groups.  For a desired tuple $t_e\in W^a$,
set the class at each root to zero.  After a parent's class
$u_p$ is fixed, set
\[
                u_v=\pi_v^*(t_e+j_{e,p}^*u_p).
\]
Its restriction difference at the parent edge is $t_e$.
Its other incident edges lead only to children whose classes
have not yet been chosen.  Proceeding down the finite tree
therefore constructs $u$ with $\rho_au=t$.

The reduced total-fiber relation and self-intersection give
\[
 \langle A_au,w\rangle
      =-\sum_e\deg_k((\rho_au)_e(\rho_{n-1-a}w)_e).
 \tag{FR17.3}\label{r17fr:energy}
\]
The numerical pairings on every native $D_v$ and $E_e$ are
perfect.  Since the second restriction map is surjective
in the complementary degree, $A_au=0$ forces
$\rho_au=0$.  The converse follows from the same identity.
This is ordinary rational perfect-pairing linear algebra,
with no higher-codimension positivity or Lefschetz premise.
\end{proof}

\begin{proposition}[Recursive Chow alignment without a comparison theorem]
\label{r17fr:Chow}
For every $v=(v_w)\in\ker A_a$, $0\le a\le n-1$, there
are actual rational Chow representatives
$\zeta_w\in\operatorname{CH}^a(D_w)_{\mathbf Q}$ of the
$v_w$ whose restrictions agree on each double stratum.
\end{proposition}

\begin{proof}
Choose arbitrary rational Chow representatives $\alpha_w$.
Keep $\zeta_w=\alpha_w$ at the roots.  Suppose the representative
$\zeta_p$ of a parent has already been fixed.  At a child $w$
with parent edge $e$, put
\[
 \Delta_e=j_{e,w}^*\alpha_w-j_{e,p}^*\zeta_p,
 \qquad
 \zeta_w=\alpha_w-\pi_w^*\Delta_e.
 \tag{FR17.4}\label{r17fr:alignment}
\]
By Lemma~\ref{r17fr:rho-surjective}, the original numerical
tuple is compatible.  Inductively $\zeta_p$ has numerical
class $v_p$, so $\Delta_e$ is numerically zero on $E_e$.
Its pullback is numerically zero on $D_w$, and $\zeta_w$
still represents $v_w$.  The right inverse identity makes
its restriction on the parent edge exactly that of
$\zeta_p$ in the Chow group.

Changing $\alpha_w$ may change restrictions at its child
edges.  No earlier equality is affected, because the only
previously treated edge at $w$ is its parent edge.  The
changes at child edges are corrected in subsequent steps.
After finitely many steps all equalities hold.  This is the
reason a forest, and not independent one-sided surjectivity
on an arbitrary graph, suffices.  The equality is already
in rational Chow groups; no assertion about a cycle-class
map from numerical groups is involved.
\end{proof}

\begin{theorem}[All-codimension strong G from parent retractions]
\label{r17fr:actual}
Under the stated model and parent-retraction hypotheses,
every geometrically homologically trivial rational Chow
class on $X$ has a numerically admissible extension on
$\mathcal X$ in every valid positive codimension.
Consequently its actual obstruction is zero, and every
specified cycle in it has a strong ordinary actual Green
form on the exact specified support complement, tested
against all compact actual forms over the original field.
For FD17's notation on this model,
\[
                   K_{\rm comp}=\ker A_{n-r}
                        \quad(1\le r\le n).
 \tag{FR17.5}\label{r17fr:compatible}
\]
\end{theorem}

\begin{proof}
Apply Proposition~\ref{r17fr:Chow} with $a=n-r$ and then
apply geometric cycle classes.  It gives exactly the
compatible representatives in KC17, with all special-fiber
multiplicities one.  KC17's arbitrary-degree compatible
lift, native top trace and rational numerical duality
produce a vertical Chow correction on this model.
The frozen literal numerical detector and exact-support
actual reconstruction give the conclusion, as in FH17.
All native residue degrees remain in the pairings.  The
theorem can likewise be used on a supplied projective
dominant witness after a finite auxiliary extension;
dominant and full-degree finite Bott--Chern injection
then descend the obstruction to the original field.
\end{proof}

\begin{corollary}[Arbitrarily twisted ruled components on a forest]
If every nonroot component is a projective-line bundle
over its parent double stratum and that stratum is one
of its sections, the preceding theorem applies.
The line-bundle twists may be arbitrary.  More generally,
any supplied smooth projective component with the stated
algebraic retraction is allowed.
\end{corollary}

\begin{proof}
The bundle projection is \eqref{r17fr:retraction}.
There is no comparison of the different section
restriction maps in this argument.  The component's
other boundary strata may therefore have distinct
Chow restriction maps.
\end{proof}

\paragraph{Relation to FT17 and the remaining graph gate.}
FT17 uses numerical divisor surjectivity together with
curve-class comparison on the doubles to obtain a
codimension-two result.  Here the algebraic retractions
supply surjectivity in all complementary degrees and
the explicit representative correction
\eqref{r17fr:alignment}.  Thus all codimensions are covered
without any numerical-to-homological assertion on those
doubles.  For a graph with cycles, the recursive procedure
returns a possible Chow discrepancy at the closing edge;
FR17 does not discard it.  FH17 treats a different sufficient
condition by common restrictions and finite actual holonomy.
No claim that arbitrary models admit either kind of data is made.
